% !PRINTED 589
\weberpage{589}
Accordingly, if $\alpha\gamma$ is divisible by $\beta$, then $\gamma$ too must be
divisible by $\beta$. We therefore have the following fundamental theorem:

5. \gesperrt{A product of two or more integers
in $R(i)$ is divisible by a prime in $R(i)$
only if at least one of its factors
is divisible by this prime.}

If one now further makes the remark that a prime can be
divisible by another only if the quotient
is a unit, that is, if the two are associated, and regards
associated primes as not essentially different, it follows:

6. \gesperrt{An integer of the field $R(i)$ can always
and in only one way be decomposed into a product
 of primes.}

For let $\alpha$ be any integer in $R(i)$. If $\alpha$ is not
itself a prime, it is decomposable, say into $\gamma\beta$; but the norms
of $\gamma$ and of $\beta$ are smaller than the norm of $\alpha$. If $\beta$
is not yet a prime, it is again decomposable, say into $\delta\varepsilon$; the
norms of these factors, which are all positive integers,
decrease continually, and one must therefore necessarily at some point come upon
a factor which is a prime. Hence every number $\alpha$
is certainly divisible by a prime $\pi$. If accordingly $\alpha = \pi\alpha'$,
then the same holds of $\alpha'$, and we have, say, $\alpha' = \pi'\alpha'' \dots$ Since the
norms of $\alpha, \alpha', \alpha'' \dots$ again all decrease, one must,
in continuing this sequence, finally come upon a unit.
Hence $\alpha$ is decomposable into a finite number of prime factors,
thus
\[
\alpha = \pi\pi'\pi'' \dots,
\]
where a unit factor remaining over at the end can be united with one of the
primes $\pi$. If now $\varkappa, \varkappa', \varkappa'' \dots$
are likewise primes and $\alpha = \varkappa\varkappa'\varkappa'' \dots$, then it follows from
\[
\varkappa\varkappa'\varkappa'' \dots = \pi\pi'\pi'' \dots
\]
by 5.\ that one of the primes $\pi$, e.g.\ the first, is divisible by $\varkappa$,
and hence is not essentially different from $\varkappa$. If one cancels
on both sides by $\varkappa = \pi$, one can repeat the same inference with $\varkappa'$,
and so forth, and thus finds that each of the primes $\varkappa$
must also occur among the $\pi$, and if among the $\varkappa$ a
prime occurs several times, it must occur at least as
often among the $\pi$. But the same inference can also be
made in the reverse direction, and from this it follows that the totality
%CONT>>
